% !PRINTED 218
\weberpage{218}
Unter den Wurzeln der Teilungsgleichung kommen auch die
$\varphi(n)$ Größen

\weq{$R_1$}{\mathrm{sn}\,h\,\Omega_1}
vor, die, wenn $h$ ein vollständiges System inkongruenter zu
$n$ teilerfremder Zahlen durchläuft, alle voneinander verschieden
sind. Das System $(R_1)$ wollen wir die \gesperrt{erste Reihe der Wurzeln}
nennen.

Ist nun $\mathrm{sn}\,\Omega_2$ eine in $(R_1)$ nicht enthaltene Wurzel, so bilden
die $\varphi(n)$ Größen

\weq{$R_2$}{\mathrm{sn}\,h\,\Omega_2,}
die sowohl untereinander als von den Wurzeln $(R_1)$ verschieden
sind, eine zweite Reihe; und auf diese Weise kann man fortfahren,
bis sämtliche $\varphi(n)\psi(n)$ Wurzeln der eigentlichen Teilungsgleichung
in $\psi(n)$ Reihen von je $\varphi(n)$ Gliedern verteilt sind.

Nach dem Multiplikationstheorem läßt sich durch eine Wurzel
jede andere Wurzel derselben Reihe \gesperrt{rational} ausdrücken in der
Form:

\weq{1}{\mathrm{sn}\,h\,\Omega = f_h(\mathrm{sn}\,\Omega),}
worin $f_h$ eine nur von dem Multiplikator $h$, nicht von der Wahl
von $\Omega$ abhängige rationale Funktion ist.

Wenn die beiden Wurzeln $x_{\mu,\mu'}$, $x_{\nu,\nu'}$ in dieselbe Reihe gehören, so muß sich die Zahl $h$ so bestimmen lassen, daß

\weq{2}{h\mu \equiv \nu, \quad h\mu' \equiv \nu' \pmod{n},}
und umgekehrt, wenn dies der Fall ist, so gehören die beiden
Wurzeln in dieselbe Reihe. Aus (2) aber folgt die Kongruenz

\weq{3}{\mu\nu' - \nu\mu' \equiv 0 \pmod{n},}
und aus dieser lassen sich auch umgekehrt die Kongruenzen (2)
wieder herleiten.

Denn da $\mu, \mu', n$ relativ prim sind, so kann man die Zahlen
$\alpha, \beta$ so bestimmen, daß
\[
\alpha\mu - \beta\mu' \equiv 1 \pmod{n}
\]
wird, und daraus folgt mittels (3):
\[
\nu \equiv (\alpha\nu - \beta\nu')\mu, \quad \nu' \equiv (\alpha\nu - \beta\nu')\mu' \pmod{n},
\]
also, wenn $\alpha\nu - \beta\nu' = h$ gesetzt wird, die Kongruenzen (2).

\gesperrt{Die Kongruenz (3) ist also die notwendige und hinreichende
Bedingung dafür, daß die beiden Wurzeln der
eigentlichen Teilungsgleichung $x_{\mu,\mu'}$, $x_{\nu,\nu'}$ derselben Reihe
angehören.}
